% Level ceiling (grade 7): v = d/t as a quotient recipe (literal
% expressions arrived in math g7); average speed = total distance over
% total time; km/h <-> m/s only via gentle whole-number examples.
% The formal quotient-quantity treatment belongs to math grade 8.
\chapter{Graphiques du mouvement et vitesse moyenne}\label{ch:g7:motion-average-speed}

Il y a deux ans, la \omterm{def:g7:motion-average-speed:formula}{vitesse} voulait dire «~des kilomètres parcourus
en une heure~», comptés sur les doigts. Maintenant, armée de lettres
et de graphiques, l'idée montre les dents~: une formule qui calcule
dans trois directions, une \omterm{prop:g5:mirrors-reflection:image}{image} qui montre tout un trajet d'un coup
d'œil --- et un piège célèbre au sujet des moyennes, qui attrape des
adultes tous les jours.

\section{La formule}

\begin{definition}[La vitesse, par la formule]\label{def:g7:motion-average-speed:formula}
Pour un trajet (ou un tronçon de trajet) parcouru à allure régulière,
la \emph{vitesse}\index{vitesse} $v$ est la distance $d$ divisée par
la \omterm{def:g2:measuring-time:duration}{durée} $t$ du parcours~:
\[
v = \frac{d}{t}.
\]
Avec $d$ en kilomètres et $t$ en \omterm{ex:g2:measuring-time:clock}{heures}, $v$ parle en
\emph{kilomètres par heure} (\unit{km/h})~; avec des \omterm{def:g2:measuring-length:units}{mètres} et des
secondes, en \emph{\omterm{def:g2:measuring-length:units}{mètres} par seconde} (\unit{m/s}) --- l'\omterm{def:g6:measurement-in-science:measuring}{unité}
préférée des scientifiques. Comme celle de la \omterm{def:g7:volume-mass-density:formula}{masse volumique}, cette
formule calcule dans les trois directions~: $v = d/t$, \
$d = v \times t$, \ $t = d/v$.
\end{definition}

\begin{example}[Les recettes au travail]\label{ex:g7:motion-average-speed:recipes}
Un train parcourt \qty{240}{km} en $2$ \omterm{ex:g2:measuring-time:clock}{heures}~:
$v = 240 \div 2 = \qty{120}{km/h}$. Quelle distance à cette allure en
$5$ \omterm{ex:g2:measuring-time:clock}{heures}~? $d = 120 \times 5 = \qty{600}{km}$. Combien de temps
pour \qty{300}{km}~? $t = 300 \div 120 = 2.5$ \omterm{ex:g2:measuring-time:clock}{heures} --- deux \omterm{ex:g2:measuring-time:clock}{heures}
et trente minutes (attention au rebelle~: $0.5$ heure fait $30$
minutes, jamais «~$50$~»).
\end{example}

\begin{example}[Deux unités, une seule vitesse]\label{ex:g7:motion-average-speed:units}
Un sprinteur court \qty{100}{m} en \qty{10}{s}~:
$v = \qty{10}{m/s}$. Cela fait combien en \omterm{def:g6:measurement-in-science:measuring}{unités} de la route~? En une
heure --- $3600$ secondes --- le sprinteur parcourrait
$10 \times 3600 = \qty{36000}{m} = \qty{36}{km}$~: donc
$\qty{10}{m/s} = \qty{36}{km/h}$. Le taux de change général~: chaque
\unit{m/s} vaut \qty{3.6}{km/h} --- parce qu'une heure contient
$3600$ secondes et un kilomètre seulement $1000$ \omterm{def:g2:measuring-length:units}{mètres}. Les
\qty{340}{m/s} du \omterm{def:g3:sound-around-us:sound}{son}, en \omterm{def:g6:measurement-in-science:measuring}{unités} de la route~: plus de
\qty{1200}{km/h}.
\end{example}

\section{Le trajet en image}

\begin{proposition}[Les lois du graphique]\label{prop:g7:motion-average-speed:graphs}
Porte la distance parcourue en fonction du temps, et le mouvement
écrit son autobiographie~:
\begin{enumerate}
  \item un \omterm{def:g6:describing-motion:uniform}{mouvement uniforme} dessine une droite --- des distances
        égales pendant des \omterm{def:g2:measuring-time:duration}{durées} égales, pas après pas~;
  \item plus la droite est \emph{raide}, plus le mouvement est
        rapide~: la pente est la \omterm{def:g7:motion-average-speed:formula}{vitesse} rendue visible~;
  \item un tronçon \emph{horizontal} est un arrêt~: le temps passe,
        la distance reste.
\end{enumerate}
Les courbures racontent les changements~: une courbe qui se redresse,
c'est une accélération~; une courbe qui s'aplatit, un ralentissement.
\end{proposition}

\begin{omfigure}
\begin{tikzpicture}
\begin{axis}[omaxis, width=10.8cm, height=6.0cm,
    xlabel={temps (\unit{min})}, ylabel={distance (\unit{km})},
    xtick={0,10,20,30,40,50,60}, ytick={0,5,10,15,20},
    ymin=0, ymax=22, xmin=0, xmax=63]
  \addplot[very thick, omDef] coordinates
    {(0,0) (20,10) (35,10) (60,20)};
  \node[right, font=\footnotesize, omDef] at (axis cs:5,4.6)
    {en route~: raide};
  \node[above, font=\footnotesize, omDef] at (axis cs:27.5,10.2)
    {horizontal~: un arrêt};
  \node[right, font=\footnotesize, omDef] at (axis cs:44,13.6)
    {de nouveau en route~: plus doux};
\end{axis}
\end{tikzpicture}

{\small La matinée d'un coursier racontée par une seule ligne~: une
course vive, un arrêt de quinze minutes, puis une allure plus douce.
Le graphique est l'autobiographie du trajet.}
\end{omfigure}

\begin{method}[Lire un graphique distance--temps]\label{met:g7:motion-average-speed:read}
\begin{enumerate}
  \item vérifier d'abord les \omterm{def:g6:sun-earth-moon:axisorbit}{axes} et leurs \omterm{def:g6:measurement-in-science:measuring}{unités} --- minutes ou
        \omterm{ex:g2:measuring-time:clock}{heures}, \omterm{def:g2:measuring-length:units}{mètres} ou kilomètres~;
  \item découper la ligne à ses cassures en tronçons~; étiqueter
        chacun~: droit et raide, droit et doux, horizontal~;
  \item pour tout tronçon droit, relever sa montée et sa longueur
        horizontale --- la distance gagnée, la \omterm{def:g2:measuring-time:duration}{durée} écoulée --- et
        diviser~: c'est la \omterm{def:g7:motion-average-speed:formula}{vitesse} du tronçon~;
  \item pour l'histoire entière, lire la distance totale au temps
        final --- et se rappeler que les tronçons horizontaux font
        partie de la \omterm{def:g2:measuring-time:duration}{durée} totale.
\end{enumerate}
\end{method}

\begin{example}[Le coursier, décodé]\label{ex:g7:motion-average-speed:decoded}
Applique la méthode à la figure. Premier tronçon~: \qty{10}{km} en
$20$ minutes --- un tiers d'heure --- donc
$v = 10 \div \tfrac{1}{3} = \qty{30}{km/h}$. Deuxième~: horizontal de
la minute $20$ à la minute $35$ --- un arrêt de livraison.
Troisième~: \qty{10}{km} en $25$ minutes, un peu moins de
\qty{24}{km/h}. Total~: \qty{20}{km} en une heure --- ce qui nous
sert sur un plateau l'idée suivante du chapitre.
\end{example}

\section{La vitesse moyenne --- et son piège célèbre}

\begin{definition}[Vitesse moyenne]\label{def:g7:motion-average-speed:average}
La \emph{vitesse moyenne}\index{vitesse moyenne} d'un trajet entier
est la distance totale divisée par la \omterm{def:g2:measuring-time:duration}{durée} \emph{totale} --- arrêts
compris~:
\[
v_{\text{moyenne}} = \frac{d_{\text{totale}}}{t_{\text{totale}}}.
\]
C'est l'unique allure régulière qui aurait couvert la même route dans
la même \omterm{def:g2:measuring-time:duration}{durée} d'ensemble --- les \qty{20}{km} du coursier en une
heure~: moyenne de \qty{20}{km/h}, alors que ses roues n'ont pas
tourné une seule fois à cette \omterm{def:g7:motion-average-speed:formula}{vitesse}.
\end{definition}

\begin{proposition}[Le piège]\label{prop:g7:motion-average-speed:trap}
La \omterm{def:g7:motion-average-speed:average}{vitesse moyenne} d'un trajet n'est \emph{pas}, en général, le
milieu de ses \omterm{def:g7:motion-average-speed:formula}{vitesses}. Les tronçons lents dévorent plus de temps que
les rapides, et pèsent donc plus lourd dans la moyenne. Seule la
recette complète --- distance totale sur \omterm{def:g2:measuring-time:duration}{durée} totale --- est digne
de confiance~; faire la moyenne des nombres de \omterm{def:g7:motion-average-speed:formula}{vitesse} eux-mêmes est
le faux pas le plus séduisant du chapitre.
\end{proposition}

\begin{example}[Le piège se referme]\label{ex:g7:motion-average-speed:sprung}
Une cycliste parcourt \qty{60}{km} à l'aller à \qty{30}{km/h}, et les
mêmes \qty{60}{km} au retour à \qty{20}{km/h}. «~Moyenne~:
\qty{25}{km/h}~»~? Vérifions honnêtement. Aller~:
$t = 60 \div 30 = 2$ \omterm{ex:g2:measuring-time:clock}{heures}. Retour~: $t = 60 \div 20 = 3$ \omterm{ex:g2:measuring-time:clock}{heures}.
Trajet entier~: \qty{120}{km} en $5$ \omterm{ex:g2:measuring-time:clock}{heures} ---
$v_{\text{moyenne}} = \qty{24}{km/h}$. La moitié lente s'est
approprié trois des cinq \omterm{ex:g2:measuring-time:clock}{heures}, et a tiré la moyenne au-dessous du
milieu. Plus tu vas vite sur une moitié, moins cette moitié existe
longtemps.
\end{example}

\begin{remark}[Ce que sait le compteur]\label{rem:g7:motion-average-speed:speedometer}
La \omterm{def:g7:motion-average-speed:average}{vitesse moyenne} décrit un trajet entier~; l'aiguille du compteur
répond à une autre question --- \emph{à quelle \omterm{def:g7:motion-average-speed:formula}{vitesse} maintenant}.
Sur le graphique, ce «~maintenant~» vit dans la pente de la ligne
\emph{en un seul point} --- facile à voir sur un tronçon droit,
subtil là où la ligne se courbe. Rendre précise la «~pente en un
point~» est l'une des plus grandes inventions des mathématiques, et
elle t'attend à la fin du lycée. En attendant~: les tronçons droits
reçoivent des nombres, les courbes reçoivent des récits.
\end{remark}

\section{Exercices}

\begin{exercise}[$\star$]\label{exo:g7:motion-average-speed:1}
Écris la formule de la \omterm{def:g7:motion-average-speed:formula}{vitesse} et ses deux \omterm{def:g6:measurement-in-science:measuring}{unités} courantes. Quelle
recette trouve une distance~? Une \omterm{def:g2:measuring-time:duration}{durée}~?
\end{exercise}

\begin{exercise}[$\star$]\label{exo:g7:motion-average-speed:2}
Un ferry parcourt \qty{45}{km} en $3$ \omterm{ex:g2:measuring-time:clock}{heures}~; un lièvre court
\qty{100}{m} en \qty{8}{s}. Calcule les deux \omterm{def:g7:motion-average-speed:formula}{vitesses}, chacune dans
son \omterm{def:g6:measurement-in-science:measuring}{unité} naturelle.
\end{exercise}

\begin{exercise}[$\star$]\label{exo:g7:motion-average-speed:3}
Convertis, avec le taux de change~: \qty{5}{m/s} en \unit{km/h}~; \
\qty{72}{km/h} en \unit{m/s}.
\end{exercise}

\begin{exercise}[$\star$]\label{exo:g7:motion-average-speed:4}
Énonce les trois lois du graphique. Que raconte une courbe qui
s'aplatit progressivement~?
\end{exercise}

\begin{exercise}[$\star$]\label{exo:g7:motion-average-speed:5}
Sur le graphique du coursier~: entre quelles minutes l'allure
est-elle la plus douce (sans être nulle)~? Comment le voir sans
calculer~?
\end{exercise}

\begin{exercise}[$\star$]\label{exo:g7:motion-average-speed:6}
Le graphique d'un marcheur montre une droite passant par les points
($30$ min, \qty{2}{km}) et ($60$ min, \qty{4}{km}). \omterm{def:g6:describing-motion:uniform}{Mouvement
uniforme} ou varié~? \omterm{def:g7:motion-average-speed:formula}{Vitesse} en \unit{km/h}~?
\end{exercise}

\begin{exercise}[$\star$]\label{exo:g7:motion-average-speed:7}
Définis la \omterm{def:g7:motion-average-speed:average}{vitesse moyenne}. Une randonnée~: \qty{12}{km} en $4$
\omterm{ex:g2:measuring-time:clock}{heures}, pique-nique d'une heure compris. \omterm{def:g7:motion-average-speed:average}{Vitesse moyenne} --- et
\omterm{def:g7:motion-average-speed:average}{vitesse moyenne} \emph{en marchant}~?
\end{exercise}

\begin{exercise}[$\star\star$]\label{exo:g7:motion-average-speed:8}
Pourquoi un tronçon lent est-il «~plus lourd~» dans une moyenne qu'un
tronçon rapide de même longueur~? Réponds avec le temps, pas avec des
formules.
\end{exercise}

\begin{exercise}[$\star\star$]\label{exo:g7:motion-average-speed:9}
Rejoue le piège~: \qty{30}{km} à l'aller à \qty{15}{km/h},
\qty{30}{km} au retour à \qty{30}{km/h}. Milieu prédit, moyenne
honnête --- et quelle moitié du trajet a possédé l'essentiel de
l'horloge~?
\end{exercise}

\begin{exercise}[$\star\star$]\label{exo:g7:motion-average-speed:10}
Esquisse (ou décris précisément) le graphique de ce trajet~:
uniforme à \qty{40}{km/h} pendant une demi-heure~; arrêt d'un quart
d'heure~; uniforme à \qty{60}{km/h} pendant un quart d'heure. Calcule
ensuite la \omterm{def:g7:motion-average-speed:average}{vitesse moyenne} du trajet.
\end{exercise}

\begin{exercise}[$\star\star$]\label{exo:g7:motion-average-speed:11}
Le comptage de l'orage, version améliorée~: le tonnerre arrive
\qty{6}{s} après l'éclair. Avec le \omterm{def:g3:sound-around-us:sound}{son} à \qty{340}{m/s}, utilise
$d = v \times t$ pour la distance de l'orage --- et confronte la
vieille règle «~diviser par trois pour avoir des kilomètres~» à ta
formule.
\end{exercise}

\begin{exercise}[$\star\star\star$]\label{exo:g7:motion-average-speed:12}
Un vieux classique qui vaut chaque minute~: un \omterm{def:g4:conductors-insulators:conductor}{conducteur} parcourt la
première moitié \emph{de la distance} d'un trajet à \qty{30}{km/h} et
veut une moyenne d'ensemble de \qty{60}{km/h}. Montre --- avec la
recette distance totale sur \omterm{def:g2:measuring-time:duration}{durée} totale, sur un trajet de
\qty{60}{km} --- que la seconde moitié devrait être parcourue en un
temps nul~: le souhait est impossible, et non pas simplement
difficile.
\end{exercise}

\section{Problème~: le vendredi du coursier}

\begin{problem}[{Problème du week-end --- une coursière à vélo, un
vendredi en ville~; le graphique du régulateur, le piège de la fiche
de prime}]\label{pb:g7:motion-average-speed:1}
Le régulateur enregistre le vendredi de la coursière Lena sur un
graphique distance--temps et en déduit sa paie. Tu vérifies la
fiche.

\medskip\noindent\textbf{Partie I --- la matinée, d'après le
graphique.}
Le graphique montre~: une montée droite de ($0$ min, \qty{0}{km})
à ($30$ min, \qty{9}{km})~; un palier jusqu'à $45$ min~; une droite
de là jusqu'à ($75$ min, \qty{15}{km})~; un palier jusqu'à $90$ min.
\begin{enumerate}
  \item Raconte l'histoire de la matinée avec des mots~: les
        tronçons, les arrêts, et ce qu'étaient vraisemblablement
        ces arrêts.
  \item \omterm{def:g7:motion-average-speed:formula}{Vitesse} sur le premier tronçon, en \unit{km/h}~?
  \item \omterm{def:g7:motion-average-speed:formula}{Vitesse} sur le second tronçon roulé~?
  \item \omterm{def:g7:motion-average-speed:formula}{Vitesse} \emph{moyenne} de Lena sur toute la matinée de $90$
        minutes, arrêts compris~?
\end{enumerate}

\medskip\noindent\textbf{Partie II --- l'après-midi, d'après la
formule.}
\begin{enumerate}[resume]
  \item Première étape de l'après-midi~: \qty{12}{km} à un régulier
        \qty{24}{km/h}. Combien de temps cela a-t-il pris, en
        minutes~?
  \item Deuxième étape~: une livraison dans le haut de la ville,
        \qty{20}{minutes} à \qty{18}{km/h}. Combien de kilomètres~?
  \item Troisième étape~: le long trajet jusqu'au dépôt,
        \qty{10}{km}, effectué en $25$ minutes. \omterm{def:g7:motion-average-speed:formula}{Vitesse} à allure
        régulière, en \unit{km/h}~?
  \item Total de l'après-midi~: additionne les distances et les
        \omterm{def:g2:measuring-time:duration}{durées} (les étapes seules, sans les pauses)~: quelle
        \omterm{def:g7:motion-average-speed:average}{vitesse moyenne} les roues ont-elles tenue en roulant~?
\end{enumerate}

\medskip\noindent\textbf{Partie III --- le piège de la fiche de
prime.}
La règle de la prime~: «~une \omterm{def:g7:motion-average-speed:average}{vitesse moyenne} sur une tournée
complète supérieure à \qty{20}{km/h} donne droit à la prime de
rapidité~».
\begin{enumerate}[resume]
  \item Le vendredi complet de Lena~: \qty{37}{km} en $3$ \omterm{ex:g2:measuring-time:clock}{heures}
        (tous arrêts compris). A-t-elle droit à la prime~?
  \item Son collègue Max la réclame~: «~J'ai roulé la moitié de ma
        tournée à \qty{30}{km/h} et l'autre moitié à
        \qty{15}{km/h} --- cela fait une moyenne de $22.5$~!~» Le
        régulateur vérifie~: les deux moitiés faisaient
        \qty{12}{km}. Calcule la vraie moyenne de Max et statue sur
        sa prime.
  \item Explique à Max, en langage de temps, où son $22.5$ a
        déraillé.
  \item Lena propose une règle de prime plus juste pour les
        coursiers --- une règle qui ne punisse pas les arrêts de
        livraison. Propose-en une (le graphique sait comment), et
        dis quelle \omterm{def:g7:motion-average-speed:formula}{vitesse} elle mesure.
\end{enumerate}

\medskip\noindent\textbf{Partie IV --- la leçon du régulateur.}
\begin{enumerate}[resume]
  \item Bilan du vendredi~: écris les trois phrases du régulateur
        --- ce que signifient sur le graphique un tronçon droit, un
        tronçon horizontal et un tronçon plus raide --- et l'unique
        avertissement au sujet des moyennes de \omterm{def:g7:motion-average-speed:formula}{vitesses}.
\end{enumerate}
\end{problem}
